\documentclass[11pt]{amsart}

\usepackage[T1]{fontenc}
\usepackage{lmodern}
\usepackage{microtype}
\usepackage{mathtools,amssymb,amsthm}
\usepackage{array}
\usepackage[hidelinks]{hyperref}

\hypersetup{
  pdftitle={An Endpoint Reveal Refutes the Exact Two-Liminal Burning Formula for Paths}
}

\newtheorem{theorem}{Theorem}
\newtheorem{proposition}[theorem]{Proposition}
\newtheorem{lemma}[theorem]{Lemma}
\theoremstyle{remark}
\newtheorem*{remark}{Remark}

\newcommand{\CL}{\operatorname{CL}}

\title[An endpoint reveal refutes an exact liminal burning formula]
{An Endpoint Reveal Refutes the Exact Two-Liminal Burning Formula for Paths}

\date{}

\subjclass[2020]{05C57, 91A43, 05C38}
\keywords{liminal burning number, graph burning, cooling number, path,
combinatorial game, counterexample}

\begin{document}

\begin{abstract}
The $k$-liminal burning number $b_k(G)$ of Bonato, Marbach, Marcoux and Mishura
interpolates between the cooling number $\CL(G)=b_1(G)$ and the burning number
$b(G)=b_{|V(G)|}(G)$: each round a saboteur reveals $k$ vertices and an
arsonist burns one revealed vertex, and the saboteur maximises the number of
rounds. A recent preprint asserts the exact formula
$b_2(P_n)=\lceil (n+2)/3\rceil$ for every $n\ge1$. We refute it. The witness is
a single two-element set: if the saboteur reveals the two \emph{endpoints} of
$P_4$ in the first round, then three rounds are needed, so $b_2(P_4)=3>2$. The
argument is four lines and needs no computation, and the lower bound
$b_2(P_4)\ge3$, which is all the refutation uses, is insensitive to the two
clauses of the definition that the sources leave ambiguous and to whether the
saboteur must reveal exactly $k$ vertices in later rounds.
\end{abstract}

\maketitle

\section{The parameter and the claim}

Let $G$ be a finite simple graph and let $k\ge1$. In the \emph{$k$-liminal
burning game} of Bonato, Marbach, Marcoux and Mishura \cite{BMMM}, two players
alternate on $G$. In round~$1$ the saboteur reveals a set of $k$ vertices and
the arsonist burns one revealed vertex, the \emph{source}. In each round
$i\ge2$, first every unburned vertex with a burned neighbour burns
(\emph{propagation}), then the saboteur reveals $k$ further vertices, then the
arsonist burns one revealed unburned vertex if one exists. Revelations are
permanent. The game ends in the round in which the last vertex burns; the
saboteur maximises the number of rounds and the arsonist minimises it, and
$b_k(G)$ is the value under optimal play. Thus $b_1(G)=\CL(G)$, the cooling
number, and $b_{|V(G)|}(G)=b(G)$, the burning number. A round in which the last
vertex burns by propagation alone is counted; this is forced by the value
$b_2(P_2)=2$ implied by the formula below.

Two clauses of the definition are read differently by the two sources.

\begin{enumerate}
\item[(A)] \emph{What may the arsonist burn from round~$2$ on?} Either any
revealed unburned vertex, or only a vertex revealed in the current round.
\item[(S)] \emph{What may the saboteur reveal?} ``$k$ additional vertices''
\cite{Target} permits re-revealing a vertex that has since burned; \cite{BMMM}
does not.
\end{enumerate}

\noindent
We fix the convention used below: from round~$2$ on the arsonist may burn any
revealed unburned vertex, and the saboteur reveals vertices that are neither
revealed nor burned. Only the upper bound in Theorem~\ref{thm:main} uses either
clause, and it uses only the first alternative in (A); the refutation itself
does not, as the theorem records.

Ambrose, Angelone, Chen, Ma, Ortiz San Miguel, Watanabe, Whitcomb and Wu
\cite{Target} state the following.

\begin{quote}
\textbf{Theorem 3.1} (\cite{Target}). For every $n\ge1$,
$b_2(P_n)=\left\lceil \dfrac{n+2}{3}\right\rceil$.
\end{quote}

\section{The counterexample}

Write $V(P_4)=\{v_1,v_2,v_3,v_4\}$ with edges $v_1v_2$, $v_2v_3$, $v_3v_4$.

\begin{theorem}\label{thm:main}
$b_2(P_4)=3$ under the convention fixed in Section~1. The lower bound
$b_2(P_4)\ge3$ holds under either alternative in (A) and in (S), and whether or
not the saboteur is required to reveal exactly two vertices in later rounds.
Since $\lceil(4+2)/3\rceil=2$, Theorem~3.1 of \cite{Target} is false.
\end{theorem}

\medskip
\noindent\textbf{The object.} The witness is the first-round reveal
\[
  S=\{v_1,v_4\},
\]
the two endpoints. Nothing else is needed; in particular no computation is
needed.

\begin{proof}
\emph{Lower bound.} Let the saboteur reveal $S=\{v_1,v_4\}$ in round~$1$. The
source must be revealed, hence $v_1$ or $v_4$; the map $v_i\mapsto v_{5-i}$ is
an automorphism of $P_4$ fixing $S$, so we may take the source to be $v_1$.
Suppose the game ended in round~$2$. At the start of round~$2$ propagation
burns exactly $v_2$, so the burned set is $N[v_1]=\{v_1,v_2\}$; the arsonist
then burns at most one further vertex. Hence at most
$|N[v_1]|+1=3$ of the four vertices are burned at the end of round~$2$, a
contradiction. Therefore $b_2(P_4)\ge3$.

This step bounds the burned set by a cardinality and never asks \emph{which}
vertex the arsonist chose, so it is valid under either alternative in (A) and in
(S) --- including the loosest, in which the arsonist may burn any unburned
vertex at all. It uses only the round-$1$ reveal, so it is also valid if the
saboteur is permitted to reveal fewer than $k$ vertices in later rounds.

\emph{Upper bound.} Let the saboteur reveal any $2$-set $S'$ in round~$1$.
If $S'$ meets $\{v_2,v_3\}$, the arsonist takes such a vertex as source; since
$N[v_2]=\{v_1,v_2,v_3\}$ and $N[v_3]=\{v_2,v_3,v_4\}$, propagation alone burns
everything by the end of round~$3$. Otherwise $S'=\{v_1,v_4\}$; the arsonist
burns $v_1$, and because revelations are permanent, $v_4$ is still revealed and
unburned in round~$2$, so the arsonist burns $v_4$ there. The burned set is
then $\{v_1,v_2,v_4\}$ and round-$3$ propagation burns $v_3$. In both cases the
game ends by round~$3$, so $b_2(P_4)\le3$. This argument uses the first
alternative in (A) and is unaffected by whether the round-$2$ reveal must have
exactly two elements or may have fewer.
\end{proof}

\begin{remark}
The claimed value $2$ requires a source $v$ with $|N[v]|+1\ge4$, i.e.\
$|N[v]|\ge3$, i.e.\ $v\in\{v_2,v_3\}$ --- precisely the two vertices the
endpoint reveal denies. Correspondingly the proof of Theorem~3.1 in
\cite{Target} has the saboteur reveal $\{v_1,v_2\}$ on turn~$1$, which hands
the arsonist the source $v_2$, and writes that it ``propose[s] a saboteur
strategy for $k=2$ on $P_n$, and claim[s] that it is optimal'' without proving
the optimality. Because the saboteur is the \emph{maximiser}, exhibiting one
saboteur strategy establishes only a lower bound on $b_2(P_n)$.
\end{remark}

\section{Scope and what was checked}

The accompanying program \verb|verify.py| mechanises the argument above on the
parsed graph: that the printed reveal is the endpoint set, that the source must
then be an endpoint, the round-$2$ cardinality bound for each possible
round-$2$ burn, and the upper bound over all six $2$-subsets of $V(P_4)$. It
also computes $b_2(P_4)=3$ by exact backward induction on the extensive-form
game under each of the four combinations of the alternatives in (A) and (S),
the saboteur there revealing exactly two vertices whenever two remain
available.

No closed form for $b_2(P_n)$ is offered, and nothing is claimed for $n\ge5$.

\begin{thebibliography}{9}

\bibitem{Target}
Ambrose, Angelone, Chen, Ma, Ortiz San Miguel, Watanabe, Whitcomb, and Wu,
\emph{Burning games on strong path products},
\href{https://arxiv.org/abs/2509.20572v2}{arXiv:2509.20572v2}, 2025.

\bibitem{BMMM}
Bonato, Marbach, Marcoux, and Mishura,
\emph{Between burning and cooling: liminal burning on graphs},
Discrete Math. \textbf{349} (2026), 115198.
\href{https://doi.org/10.1016/j.disc.2026.115198}{doi:10.1016/j.disc.2026.115198};
\href{https://arxiv.org/abs/2505.10727}{arXiv:2505.10727}.

\end{thebibliography}

\end{document}
